\documentclass{beamer}
\usetheme{Berkeley}
\usepackage[utf8]{inputenc}
\usepackage[brazil]{babel}
\usepackage{graphicx}
\usepackage{tcolorbox}
\let\raggedright\justifying

\tcbset{
  roteiro/.style={
    colback=yellow!10,
    colframe=orange!70!black,
    fonttitle=\bfseries,
    title=Roteiro do Apresentador,
    boxrule=0.8pt,
    rounded corners,
    before skip=4pt,
    after skip=4pt
  }
}

\title{Análise Estatística para Agricultura}
\author{
    Grupo de Apresentadores: Beatriz, Thalita, Jamile, Sayore, Vitor, Hanna, Asary, Daniel, Alberto e Felipe
}
\date{}

\begin{document}

% ==========================================================
% ====================== CAPA ==============================
% ==========================================================


\frame{
    \titlepage
    \vspace{0.5cm}
    \begin{tcolorbox}[colback=yellow!10,colframe=orange!70!black,boxrule=0.8pt,rounded corners]
        \textbf{Observações para a apresentação:} \\
        Cada apresentador fala apenas 15–20 segundos. \\
        Apenas pontos-chave, sem detalhes longos.
    \end{tcolorbox}
}

% ==========================================================
% ====================== APRESENTADOR 1 ====================
% ==========================================================


\begin{frame}{Beatriz — Contextualização}
\begin{itemize}
    \item Projeto para monitoramento ambiental agrícola.
    \item Coleta de dados de temperatura, umidade do ar e solo.
\end{itemize}

\begin{tcolorbox}[roteiro]
Falar de forma objetiva sobre o projeto e sua relevância.
\end{tcolorbox}
\end{frame}

% ==========================================================
% ====================== APRESENTADOR 2 ====================
% ==========================================================

\begin{frame}{Thalita — Desafios}
\begin{itemize}
    \item Variabilidade climática e do solo.
    \item Necessidade de coleta confiável e contínua.
\end{itemize}

\begin{tcolorbox}[roteiro]
Apresentar rapidamente os principais desafios do projeto.
\end{tcolorbox}
\end{frame}

% ==========================================================
% ====================== APRESENTADOR 3 ====================
% ==========================================================

\begin{frame}{Jamile — Objetivos}
\begin{itemize}
    \item Monitorar variáveis ambientais.
    \item Apoiar decisões agrícolas com base em dados.
\end{itemize}

\begin{tcolorbox}[roteiro]
Destacar os objetivos gerais e específicos em poucas palavras.
\end{tcolorbox}
\end{frame}

% ==========================================================
% ====================== APRESENTADOR 4 ====================
% ==========================================================

\begin{frame}{Sayore — Justificativa}
\begin{itemize}
    \item Dados precisos reduzem desperdícios.
    \item Integração entre tecnologia e sustentabilidade.
\end{itemize}

\begin{tcolorbox}[roteiro]
Ressaltar relevância do projeto de forma concisa.
\end{tcolorbox}
\end{frame}

% ==========================================================
% ====================== APRESENTADOR 5 ====================
% ==========================================================

\begin{frame}{Vitor — Métodos: Coleta}
\begin{itemize}
    \item Arduino Uno com sensores DHT11 e HD-38.
    \item Coleta automática de dados em campo.
\end{itemize}

\begin{tcolorbox}[roteiro]
Mostrar sensores rapidamente, apenas citar funções.
\end{tcolorbox}
\end{frame}

% ==========================================================
% ====================== APRESENTADOR 6 ====================
% ==========================================================

\begin{frame}{Hanna — Métodos: Processamento}
\begin{itemize}
    \item Python: Pandas, NumPy, Matplotlib.
    \item Análise rápida de padrões e tendências.
\end{itemize}

\begin{tcolorbox}[roteiro]
Explicar o processamento de forma objetiva, sem detalhes técnicos longos.
\end{tcolorbox}
\end{frame}

% ==========================================================
% ====================== APRESENTADOR 7 ====================
% ==========================================================

\begin{frame}{Asary — Resultados Principais}
\begin{itemize}
    \item Temperatura e umidade estáveis.
    \item Solo com alta retenção hídrica.
\end{itemize}

\begin{tcolorbox}[roteiro]
Apontar rapidamente métricas-chave com gráficos ou imagens de apoio.
\end{tcolorbox}
\end{frame}

% ==========================================================
% ====================== APRESENTADOR 8 ====================
% ==========================================================

\begin{frame}{Daniel — Entregáveis}
\begin{itemize}
    \item Relatórios parciais conforme plano.
    \item Evidências objetivas do avanço do projeto.
\end{itemize}

\begin{tcolorbox}[roteiro]
Citar entregáveis de forma direta.
\end{tcolorbox}
\end{frame}

% ==========================================================
% ====================== APRESENTADOR 9 ====================
% ==========================================================

\begin{frame}{Alberto — Engajamento}
\begin{itemize}
    \item Participação de alunos e parceiros.
    \item Exemplos práticos de colaboração.
\end{itemize}

\begin{tcolorbox}[roteiro]
Apresentar engajamento de forma objetiva, com fotos ou casos rápidos.
\end{tcolorbox}
\end{frame}

% ==========================================================
% ====================== APRESENTADOR 10 ===================
% ==========================================================

\begin{frame}{Felipe — Próximos Passos}
\begin{itemize}
    \item Conclusão da análise completa.
    \item Recomendações para manejo sustentável.
    \item Divulgação e publicações futuras.
\end{itemize}

\begin{tcolorbox}[roteiro]
Finalizar a apresentação, agradecer e abrir espaço para perguntas.
\end{tcolorbox}
\end{frame}

\end{document}

